\documentclass[8pt]{strategynet-journal}
\usepackage{volt-paper-style}
\hypersetup{pdftitle={Live Option Analytics with Volt and Hypercube},pdfauthor={research@strategynet.ai}}
\newcommand{\VOLT}{\textsf{Volt}}
\newcommand{\HYPER}{\textsf{Hypercube}}
\SNHeader{strategynet.ai:Volt + Hypercube}
\title{Live Option Analytics with \VOLT{} and \HYPER{}\\
\large Prices, Volatility Surfaces and Portfolio Exposures}
\author{\href{mailto:research@strategynet.ai}{research@strategynet.ai}}
\date{October 2026}
\begin{document}
\maketitle
\begin{abstract}
\normalsize An options trader sees prices moving across hundreds of strikes and expiries while the underlying stock moves alongside them. Which changes come from the stock, which reflect a change in the market's assessment of volatility, and how much of either movement reaches the portfolio? We develop the calculations behind these questions using a European call quoted near 6.09 and a position of ten short contracts. The example connects bid and ask prices to implied volatility, builds the coordinates of a volatility surface, and shows how a stock hedge leaves substantial curvature and volatility exposure. We then extend the calculation across stocks: equity-factor loadings translate option delta into common return exposures, while separate volatility loadings describe changes in the surfaces themselves. This provides a practical basis for comparing securities, locating concentrated risks and investigating unusual option prices. Volt supplies the numerical calculations and Hypercube organizes the observations and results. The implemented single-underlying application provides the starting point for the cross-sectional analysis developed here.
\end{abstract}
\SNtwocolumntableofcontents
\begin{SNcolumns}
\section{Reading the option market}
A stock rises by one dollar and several calls rise with it. Most of that move
may be explained by their exposure to the stock. Elsewhere in the chain, put
prices rise even though the stock has barely moved. Perhaps investors are
paying more for downside protection, an earnings announcement is approaching,
or the quoted spread has widened. A useful options view helps the reader work
through these possibilities, from the observed prices to their effect on a
position.

A ten-dollar option price can reflect a valuable exercise right already deep in the money, a long
remaining life, or substantial uncertainty about the underlying. Implied
volatility expresses a price through a specified option model and helps compare
strikes and maturities. Greeks describe how the model value responds to changes
in its inputs. Together, these calculations let us ask whether a price change
follows the stock, shifts the volatility smile, or changes the risk of a book.

Consider a European call with stock price $S=100$, strike $K=100$ and six months
remaining, $T=0.5$ years. The annual continuously compounded interest rate is
$r=3\%$ and the continuous dividend yield is $q=1\%$. Its bid is 6.00 and its ask
is 6.18 per underlying unit. With a contract multiplier of 100, the midpoint
6.09 corresponds to a contract value of 609. A trade at the ask would cost 618,
before fees. The distinction matters whenever we compare a theoretical mark
with an executable price.

Our examples use European exercise, a continuous dividend yield and the
Black--Scholes--Merton pricing convention~\cite{bs,merton}. American rights
allow exercise before expiry and require an exercise-aware valuation engine. Across a chain, the stock observation and carry assumptions also
need to refer to a common valuation time. Otherwise an apparent change in
volatility can simply reflect an option quote paired with a newer stock price.

\section{What a quote implies about volatility}
For the call above, the model price is
\begin{align}
 C&=Se^{-qT}N(d_1)-Ke^{-rT}N(d_2),\\
 d_1&=\frac{\log(S/K)+(r-q+\tfrac12\sigma^2)T}{\sigma\sqrt T},
 \qquad d_2=d_1-\sigma\sqrt T.
\end{align}
Here $N$ is the standard normal distribution function, and $\sigma$ is annual
volatility expressed as a fraction. Entering $\sigma=0.20$ gives a call value of
6.090127, very close to the observed midpoint. Implied volatility reverses this
calculation: we supply the price and solve for $\sigma$.

A numerical solver can use Newton's method,
\[
 \sigma_{j+1}=\sigma_j-
 \frac{C(\sigma_j)-C_{\rm mid}}{C_\sigma(\sigma_j)},
\]
where $C_\sigma$ is vega, the price sensitivity to volatility. Starting at 25\%
for this quote gives 19.997896\% after one update and 19.999542\% after the next.
Repricing at the latter value recovers the midpoint 6.09. Volt combines these
updates with a bracketed search, keeping the candidate inside a valid range.
The European call price limits supply an initial financial check; in this
example they run from 0.990054 to 99.501248.

The bid of 6.00 and ask of 6.18 imply the volatility interval in
Figure~\ref{fig:ivspread}. A trader
buying at the ask pays for roughly 20.3234\% volatility under the stated inputs;
selling at the bid realizes roughly 19.6757\%.

\begin{center}\begin{minipage}{\linewidth}\centering
\begin{tikzpicture}[
 box/.style={draw=darkblue,rounded corners,align=center,
 text width=1.56cm,minimum height=.91cm,font=\normalsize},
 arrow/.style={-{Latex[length=1.5mm]},draw=darkblue}]
\node[box] (bid) {Bid 6.00\\IV 19.6757\%};
\node[box,right=.27cm of bid] (mid) {Mid 6.09\\IV 19.9995\%};
\node[box,right=.27cm of mid] (ask) {Ask 6.18\\IV 20.3234\%};
\draw[arrow] (bid)--(mid);\draw[arrow] (mid)--(ask);
\end{tikzpicture}
\captionof{figure}{Bid, midpoint and ask translated into annual implied volatility for the same six-month call.}\label{fig:ivspread}
\end{minipage}\end{center}

The local sensitivity of implied volatility to a price change is the reciprocal
of vega:
\[
 d\sigma_{\rm imp}\simeq\frac{dC}{C_\sigma}.
\]
The six-month call has vega 27.7895 per unit of volatility, so a one-cent price
change implies about 0.036 volatility percentage points. With only two calendar
days remaining, an otherwise comparable at-the-money call has vega about
2.9526; the same cent implies 0.339 points. At vega 0.25 it implies four points.
A striking IV move in a very small premium can therefore arise from a single
price tick. Reading it alongside the premium, spread and time to expiry gives
a much clearer account of the observation.

A narrow, recent two-sided market anchors a fit more strongly than a wide
or old quote when observations are weighted by their price uncertainty. Comparing fitted prices with the
original bid and ask keeps the assessment in units a trader can act on.

\section{Surfaces and price sensitivities}
An option chain gives a set of implied volatilities at different strikes and
expiries. A volatility surface connects them. To compare strikes across
maturities we use the forward $F_T=S e^{(r-q)T}$ and log-forward moneyness
$k=\log(K/F_T)$. Total variance is $w(k,T)=\sigma(k,T)^2T$. For our call,
\[
 F_T=101.005017,\qquad k=-0.010000,\qquad w=0.020000.
\]
These coordinates separate a strike's position relative to the forward from
the variance accumulated over the remaining life. For example, annual
volatilities of 21\%, 20\% and 19\% at three, six and twelve months correspond
to total variances of 0.011025, 0.020000 and 0.036100. The downward-sloping
volatility term structure still implies more accumulated variance at longer
horizons.

The current Volt surface interpolates total variance linearly between
log-moneyness points within an expiry and then between expiries. It answers
queries within the supported strike and maturity range. This gives the
application a transparent way to compare quotes and value intermediate
contracts. Financial consistency adds another requirement. At a fixed expiry,
call prices should decline with strike and form a convex curve. Put--call
parity links calls and puts to the forward. The implementation checks these
relations on its inspection grid, together with calendar total-variance
consistency. Constrained surface models, including the arbitrage-free SVI
construction of Gatheral and Jacquier~\cite{svi}, provide a stronger framework
for controlling the shape throughout the fitted domain.

The second strike derivative of a smooth European call-price curve recovers
a risk-neutral density,
\[
 f_Q(K)=e^{rT}\frac{\partial^2 C}{\partial K^2},
\]
as shown by Breeden and Litzenberger~\cite{bl}. A negative value points to
inconsistent option prices or an unsuitable fit. A smooth-looking picture can
therefore still require substantial correction before it supports valuation.

For a contract valued at a specified volatility, the familiar sensitivities are
\begin{align*}
 \Delta=C_S&=e^{-qT}N(d_1),\\
 \Gamma=C_{SS}&=\frac{e^{-qT}\phi(d_1)}{S\sigma\sqrt T},\\
 \nu=C_\sigma&=Se^{-qT}\phi(d_1)\sqrt T.
\end{align*}
Here $\phi$ is the standard normal density. Our call has delta 0.553457, gamma 0.0277895 and vega 27.789495. Delta measures
the local value change for a dollar move in stock while holding volatility
fixed. Gamma measures how that delta changes. A one-percentage-point volatility
rise, from 20\% to 21\%, uses $0.01\nu$, approximately 0.277895 per underlying
unit.

Suppose the smile is held fixed in log-forward moneyness, with local
volatility slope $\sigma_k=-0.30$. When spot rises, our fixed-strike call
moves along that smile. Since $dk/dS=-1/S$, its volatility changes at
$\sigma_S=0.003$ per dollar, giving
\[
 \frac{dC}{dS}=C_S+C_\sigma\sigma_S=0.636826.
\]
The difference from 0.553457 comes from the accompanying volatility movement.
The application supports fixed-volatility partials and spot shocks that retain
the fitted log-forward smile. An empirical model of smile dynamics would
supply a further, data-dependent scenario rule.

If we change spot and re-solve volatility while requiring the call price
to remain exactly 6.090127, the
price stays at its target by construction. Differentiating that calibration
gives $C_S+C_\sigma\sigma_S=0$. For the example, the required volatility
adjustment is about $-0.019916$ per dollar. This calculation describes how the
fitted parameter compensates for a changed input. For hedge risk, we instead
specify the joint market movement we wish to evaluate and reprice along it.
Automatic differentiation follows whichever dependencies we provide.

\section{A stock hedge and the risks that remain}
Now take ten short calls, each with multiplier 100. Signed contract quantities
turn the per-unit calculations into portfolio amounts:
\[
 V=\sum_i n_i m_i C_i,\qquad
 G_V=\sum_i n_i m_i G_i.
\]
Here $n_i$ is signed contracts, $m_i$ is the multiplier and $G_i$ is a compatible
per-unit sensitivity. The short calls have the following exposures at the
unrounded theoretical mark.
\begin{center}\normalsize
\begin{tabular}{lr}
\toprule
Quantity & Short-call position\\
\midrule
Option value&$-6{,}090.127$\\
Delta, stock units&$-553.457$\\
Gamma, currency / squared price unit&$-27.7895$\\
Vega, currency per volatility point&$-277.895$\\
Theta, currency per calendar day&$+17.7592$\\
Rho, currency per rate basis point&$-2.46278$\\
\bottomrule
\end{tabular}
\end{center}
Theta measures calendar decay, using $-C_T/365$ here. A rate basis point is
0.0001 in the annual rate. The option value belongs alongside the premium cash,
stock position and financing account when calculating the complete book.

Buying 553.457 stock units offsets the initial delta. With whole shares, buying
553 leaves a small residual delta of approximately $-0.457$. Even the exact
fractional hedge retains curvature and volatility exposure, as the two
instantaneous shocks in Figure~\ref{fig:hedgeshocks} show.

\begin{center}\begin{minipage}{\linewidth}\centering
\begin{tikzpicture}[box/.style={draw=darkblue,rounded corners,align=center,
 font=\normalsize,text width=2.45cm,minimum height=.83cm},
 arrow/.style={-{Latex[length=1.5mm]},draw=darkblue}]
\node[box,text width=5.55cm] (book) {Ten short calls + 553.457 stock units\\Initial fixed-volatility delta: zero};
\node[box,below=.55cm of book,xshift=-1.6cm] (spot) {Stock: 100 to 101\\Book change: $-13.80$};
\node[box,below=.55cm of book,xshift=1.6cm] (vol) {Volatility: 20\% to 21\%\\Book change: $-277.89$};
\draw[arrow] (book.south)--(spot.north);\draw[arrow] (book.south)--(vol.north);
\end{tikzpicture}
\captionof{figure}{Full repricing of two separate shocks, with time and carry fixed. The stock shock holds volatility fixed; the volatility shock holds stock fixed.}\label{fig:hedgeshocks}
\end{minipage}\end{center}

For the stock move, the call reprices to 6.657382. The short calls lose 567.254
and the stock gains 553.457, giving a combined loss of 13.797. The local gamma
estimate is $\tfrac12\Gamma_V(\Delta S)^2=-13.895$, close to full repricing for
this move. For the volatility shock, the call reprices to 6.368020 and the short
calls lose 277.893. The stock value stays unchanged. A headline delta of zero
therefore leaves economically significant exposures to monitor.

Under the log-forward smile rule above, stock 101 implies volatility
$0.20+0.30\log(101/100)=0.2029851$. The call then reprices to 6.740140. Comparing
this result with the fixed-volatility result makes the consequence of the
assumed smile movement concrete.

For a larger book, expiry and strike buckets locate the source of aggregate
vega. Their calculation follows the surface mapping. If an option's volatility
is a linear blend of two bucket levels with weights 0.6 and 0.4, its bucket
vegas are $0.6\nu$ and $0.4\nu$. Interpolation in total variance produces different
weights through the chain rule. Full scenario repricing then checks the local
approximation, especially for large moves, short expiries and concentrated
positions. Realized trading results additionally include changes in holdings,
execution prices, fees, dividends and funding.

\section{Across stocks, sectors and factors}
Six-month at-the-money volatility,
a fixed-moneyness downside skew and the difference between short- and
long-dated volatility each describe a specific part of a stock's surface.
Displaying their changes across a sector can reveal whether a repricing is
broad or concentrated in one name. A useful exception might be an increase in
one stock's downside volatility while its peers remain stable. The next step
is to inspect its quotes, upcoming events and liquidity, then compare the
change with that stock's own history. A wide spread can explain an apparent
outlier before an economic explanation is needed.

Equity-factor exposures translate stock moves into a common portfolio scale.
Let the simple stock return satisfy $r_i=\sum_a B_{ia}f_a$, where $f_a$ is a
factor return and $B_{ia}$ is the stock's estimated loading. If $\Delta_i$ includes
both option and stock holdings, define cash delta $D_i=S_i\Delta_i$. Then
\[
 \Delta V\simeq\sum_a E_a f_a,\qquad
 E_a=\sum_i D_i B_{ia}.
\]
For cash deltas of $+5{,}000$ and $-3{,}000$, and loadings 1.2 and 0.8 to the same
factor, $E=3{,}600$. A 1\% factor rise implies a first-order gain of 36. This
comparison gives stock, sector and common-factor views of the same book. Shared
currency units and the stated simple-return mapping make the aggregation
interpretable; curvature and accompanying volatility movements enter fuller
scenarios.

The equity loadings above map factor returns into stock-price changes.
For vega exposure, we need loadings that map volatility-factor changes into
movements of the option surface. Suppose
$\delta\sigma_i=\sum_b L_{ib}z_b$. With vega $\nu_i$ measured per unit of annual
volatility, the corresponding exposure is
\[
 H_b=\sum_i n_i m_i\nu_i L_{ib}.
\]
The shocks $z_b$ might represent a common volatility rise, a steeper downside
smile or a change in the term structure. Estimating these movements requires a
history of surfaces on comparable coordinates. It also allows a portfolio to
carry little net equity-factor delta while retaining a large common volatility
exposure.

Historical changes in surface levels, skews and term slopes can supply
option-market factors. After scaling those changes consistently, principal
components extract the combinations accounting for the most shared variation.
An alternative first explains the changes with existing equity and sector
factors, then studies the remaining common variation. This second route
asks what information the option market adds to the return-factor model.
Raw prices and their Greeks have strong mathematical relationships; choosing
financially comparable features and evaluating them on later observations
helps distinguish a useful new factor from a restatement of existing inputs.

A cross-sectional map
shows comparable surface changes by stock and sector. An exposure view shows
the signed book's equity- and volatility-factor sensitivities. An exception
list brings together unusual peer-relative changes and the quotes underlying
them. Selecting a stock can connect all three views to its chain and to a
repriced scenario. Price-change attribution adds another useful link:
\[
 \Delta C\simeq\Delta\,\Delta S+
 \tfrac12\Gamma(\Delta S)^2+\nu\,\Delta\sigma+\Theta\,\Delta t.
\]
For an interval with unchanged holdings, this explains the local stock,
curvature, volatility and calendar contributions. Full repricing supplies the
comparison; the remaining difference includes higher-order and interaction
terms. Trade cash flows and changes in holdings complete an attribution of
realized portfolio P\&L.

Volt currently provides the double-precision CPU inversion, surface and risk
calculations for a single underlying~\cite{volt}. Hypercube organizes the
observations and calculated results~\cite{hypercube}. In a local scaling test on
an Apple M3 running other applications, a synthetic 1,296-quote chain took a
median 8.225 milliseconds for the numerical calculation and 15.188 milliseconds
for the complete local request and response. The test used 20 updates after three warm-up submissions. One millisecond is
one thousandth of a second. The request interval includes parsing and
serialization; feed arrival and browser rendering require separate measurements.

\end{SNcolumns}
\clearpage
\begin{SNcolumns}
\begin{thebibliography}{9}
\interlinepenalty=10000
\bibitem{bs} F. Black and M. Scholes. The pricing of options and corporate
liabilities. \emph{Journal of Political Economy} 81(3), 637--654, 1973.
\url{https://doi.org/10.1086/260062}.
\bibitem{merton} R. C. Merton. Theory of rational option pricing.
\emph{Bell Journal of Economics and Management Science} 4(1), 141--183, 1973.
\url{https://doi.org/10.2307/3003143}.
\bibitem{bl} D. T. Breeden and R. H. Litzenberger. Prices of state-contingent
claims implicit in option prices. \emph{Journal of Business} 51(4), 621--651, 1978.
\url{https://doi.org/10.1086/296025}.
\bibitem{svi} J. Gatheral and A. Jacquier. Arbitrage-free SVI volatility surfaces.
\emph{Quantitative Finance} 14(1), 59--71, 2014. Author preprint:
\url{https://arxiv.org/abs/1204.0646}.
\bibitem{hypercube} StrategyNet. \emph{Hypercube}: architecture and publication
interfaces, source documentation, October 2026.
\url{https://github.com/ingvor-restrat/hypercube}.
\bibitem{volt} StrategyNet. \emph{Volt}: numerical companion and financial
semantics, source documentation, October 2026.
\url{https://github.com/ingvor-restrat/volt}.
\end{thebibliography}
\end{SNcolumns}
\end{document}
